
\subsubsection{3.1 Dataset}

The Anthropic HH-RLHF dataset \citep{bai2022training} contains 169,352 preference pairs: 160,800 in the train split and 8,552 in the test split. Each pair includes a chosen response (preferred by human annotators) and a rejected response to the same conversational context. The matched-pair structure controls for conversation topic, user question, and dialogue history.

\subsubsection{3.2 Linguistic Marker Operationalization}

Three marker types were extracted:

\textbf{Modal Verbs (Primary):} Counts of could, might, should, may, would per 100 words, extracted via spaCy v3.5 part-of-speech tagging (TAG=='MD').

\textbf{Hedging Markers (Secondary):} Counts of appears, seems, possibly, perhaps, might, likely per 100 words via lexicon matching.

\textbf{Alternative-Framing (Secondary):} Counts of phrases including ''on the other hand,'' ''alternatively,'' ''another option,'' ''you could also'' per 100 words via regex pattern matching.

All counts were normalized to per-100-words frequency to control for response length differences.

\subsubsection{3.3 Validation Protocol}

\textbf{Stage 1 (H-E1, Existence):} Tests whether markers can be reliably extracted with sufficient distributional variance. Success criteria: coefficient of variation (CV) > 0.3 and extraction precision > 90\%.

\textbf{Stage 2 (H-M, Mechanism):} Tests construct validity through four criteria:
\begin{enumerate}
\item Effect size: Cohen's d $\geq$ 0.15 with p < 0.05 (paired t-test)
\item Internal consistency: Cronbach's $\alpha$ > 0.7 across three marker types
\item Directional hypothesis: chosen responses have fewer markers than rejected
\item Cross-split replication: effect replicates in both train and test splits
\end{enumerate}

The threshold d=0.15 was selected to represent a minimally meaningful effect size for practical deployment decisions. Cronbach's $\alpha$ > 0.7 is the standard psychometric threshold for acceptable internal consistency \citep{nunnally1978psychometric}.

\subsubsection{3.4 Statistical Analysis}

Paired t-tests were conducted on within-pair differences (chosen - rejected) for each marker type. Cohen's d for paired samples was calculated as mean(differences) / SD(differences). Cronbach's alpha was computed across the three marker types treating them as items measuring a latent construct. All analyses were performed using Python with scipy.stats and standard statistical libraries.

\subsubsection{3.5 Power Analysis}

With N=169,352 pairs, the study achieves >0.99 power to detect effects of d=0.15 at $\alpha =0.05$. The test split alone (N=8,552) provides 0.95 power for the threshold effect size.
